
\documentclass{article}
\usepackage{colortbl}
\usepackage{makecell}
\usepackage{multirow}
\usepackage{supertabular}

\begin{document}

\newcounter{utterance}

\centering \large Interaction Transcript for game `function\_detective', experiment `two\_numbers\_passive\_labeled\_pairs', episode 249 with mistral{-}large{-}3{-}azure.
\vspace{24pt}

{ \footnotesize  \setcounter{utterance}{1}
\setlength{\tabcolsep}{0pt}
\begin{supertabular}{c@{$\;$}|p{.15\linewidth}@{}p{.15\linewidth}p{.15\linewidth}p{.15\linewidth}p{.15\linewidth}p{.15\linewidth}}
    \# & \multicolumn{2}{c}{Player} && \multicolumn{2}{c}{Game Master} \\
    \hline

    \theutterance \stepcounter{utterance}  
    & & & \multicolumn{4}{p{0.6\linewidth}}{
        \cellcolor[rgb]{0.9,0.9,0.9}{
            \makecell[{{p{\linewidth}}}]{
                \texttt{\tiny{[P1$\langle$GM]}}
                \texttt{I have a hidden Python function for you to figure out.} \\
\\ 
\texttt{Your job is to infer what it does and then submit a correct Python implementation.} \\
\\ 
\texttt{Function signature:} \\
\texttt{(a: int, b: int) {-}> int} \\
\\ 
\texttt{You have 15 turns total.} \\
\\ 
\texttt{I reveal candidate input{-}output pairs one at a time and tell you whether each one is correct.} \\
\\ 
\texttt{You must make your best SOLVE attempt no later than the second{-}to{-}last turn.} \\
\texttt{Do not wait until the very end.} \\
\\ 
\texttt{You can use these actions:} \\
\texttt{{-} JSON format described at the end of the prompt} \\
\texttt{{-} SOLVE: ```python ... ```} \\
\\ 
\texttt{Do not add import statements in your guess.} \\
\texttt{Anything you might need is already available, so just write the function body in hypothesis and return the answer.} \\
\\ 
\texttt{Preloaded examples:} \\
\texttt{1. ((({-}10, 10), 0), True)} \\
\\ 
\texttt{Scoring:} \\
\texttt{{-} correctness is required for any non{-}zero score} \\
\texttt{{-} in interactive modes, solving with fewer information requests is better} \\
\\ 
\texttt{Input formatting reminders:} \\
\texttt{{-} strings must be quoted} \\
\texttt{{-} booleans must be True/False} \\
\texttt{{-} use commas between multiple values} \\
\texttt{{-} do not write any extra action tags in the same message} \\
\\ 
\texttt{Format example:} \\
\texttt{You write:} \\
\texttt{\{'current\_mode': <probing/confirming/don't know; choose whichever best describes your current state in the task>, 'input\_rationale': <brief description of why you chose this hypothesis>, 'hypothesis': <either the exact string "unknown" if you have no guess yet OR your current best guess as runnable Python matching the given signature, with a one{-}line docstring stating what it does. Never write placeholder code, if you cannot commit to a guess, use "unknown". Match the given signature.} \\
\texttt{e.g.} \\
\texttt{def f(x):} \\
\texttt{	"""Returns the absolute value of x."""} \\
\texttt{	return x if x >= 0 else {-}x} \\
\\ 
\texttt{Use "unknown" if you have no guess.>, 'confidence\_score': <score between 1{-}100. 100 means highest confidence in your hypothesis>\}} \\
\\ 
\texttt{Output only one valid json per turn.} \\
            }
        }
    }
     \\ \\

    \theutterance \stepcounter{utterance}  
    & \multicolumn{4}{p{0.6\linewidth}}{
        \cellcolor[rgb]{0.561,0.737,0.561}{
            \makecell[{{p{\linewidth}}}]{
                \texttt{\tiny{[P1$\rangle$GM]}}
                \texttt{```json} \\
\texttt{\{} \\
\texttt{    "current\_mode": "probing",} \\
\texttt{    "input\_rationale": "The preloaded example (({-}10, 10), 0) suggests that the function might return the sum of a and b, but since {-}10 + 10 = 0, this fits. However, it could also be other operations like multiplication (which would not fit) or bitwise operations. I need more examples to confirm.",} \\
\texttt{    "hypothesis": "unknown",} \\
\texttt{    "confidence\_score": 10} \\
\texttt{\}} \\
\texttt{```} \\
            }
        }
    }
    & & \\ \\

    \theutterance \stepcounter{utterance}  
    & & & \multicolumn{4}{p{0.6\linewidth}}{
        \cellcolor[rgb]{0.9,0.9,0.9}{
            \makecell[{{p{\linewidth}}}]{
                \texttt{\tiny{[P1$\langle$GM]}}
                \texttt{OUTPUT: (((2, {-}1), {-}1), False)} \\
            }
        }
    }
     \\ \\

    \theutterance \stepcounter{utterance}  
    & \multicolumn{4}{p{0.6\linewidth}}{
        \cellcolor[rgb]{0.561,0.737,0.561}{
            \makecell[{{p{\linewidth}}}]{
                \texttt{\tiny{[P1$\rangle$GM]}}
                \texttt{```json} \\
\texttt{\{} \\
\texttt{    "current\_mode": "probing",} \\
\texttt{    "input\_rationale": "The new example ((2, {-}1), {-}1) is False, meaning the function does not return {-}1 for inputs (2, {-}1). If the function were addition, 2 + ({-}1) = 1, not {-}1, so this still fits. However, if it were subtraction (a {-} b), 2 {-} ({-}1) = 3, which is not {-}1. If it were multiplication, 2 * ({-}1) = {-}2, which is also not {-}1. This rules out simple arithmetic operations. The next guess could be bitwise AND or OR, but 2 \& {-}1 = 2 (not {-}1) and 2 | {-}1 = {-}1 (which matches the candidate output but was marked False). This suggests the function is not a bitwise operation. Another possibility is the minimum or maximum of the two numbers, but min(2, {-}1) = {-}1 (which was marked False), so it's not min. Max(2, {-}1) = 2, which is not {-}1. I need to explore other operations like absolute difference or something else.",} \\
\texttt{    "hypothesis": "unknown",} \\
\texttt{    "confidence\_score": 30} \\
\texttt{\}} \\
\texttt{```} \\
            }
        }
    }
    & & \\ \\

    \theutterance \stepcounter{utterance}  
    & & & \multicolumn{4}{p{0.6\linewidth}}{
        \cellcolor[rgb]{0.9,0.9,0.9}{
            \makecell[{{p{\linewidth}}}]{
                \texttt{\tiny{[P1$\langle$GM]}}
                \texttt{OUTPUT: ((({-}10, 1000), {-}10000), True)} \\
            }
        }
    }
     \\ \\

    \theutterance \stepcounter{utterance}  
    & \multicolumn{4}{p{0.6\linewidth}}{
        \cellcolor[rgb]{0.561,0.737,0.561}{
            \makecell[{{p{\linewidth}}}]{
                \texttt{\tiny{[P1$\rangle$GM]}}
                \texttt{```json} \\
\texttt{\{} \\
\texttt{    "current\_mode": "probing",} \\
\texttt{    "input\_rationale": "The new example (({-}10, 1000), {-}10000) is True, which strongly suggests that the function is computing the product of the two inputs (since {-}10 * 1000 = {-}10000). The first example (({-}10, 10), 0) also fits this hypothesis ({-}10 * 10 = {-}100, but the output was 0, which contradicts). Wait, no: {-}10 * 10 is {-}100, not 0, so this contradicts the first example. Thus, it cannot be simple multiplication. Another possibility is that the function returns a * b if a and b have opposite signs, but the first example ({-}10, 10) would then be {-}100, not 0. Alternatively, it could be a * b only if a is negative, but the second example (2, {-}1) would then be {-}2, not {-}1. This is confusing. Maybe it's a * b when b is positive, but the second example (2, {-}1) would then be {-}2, not {-}1. Alternatively, it could be a * abs(b) or abs(a) * b, but let's test: for ({-}10, 100} \\
            }
        }
    }
    & & \\ \\

    \theutterance \stepcounter{utterance}  
    & & & \multicolumn{4}{p{0.6\linewidth}}{
        \cellcolor[rgb]{0.9,0.9,0.9}{
            \makecell[{{p{\linewidth}}}]{
                \texttt{\tiny{[P1$\langle$GM]}}
                \texttt{Game Violation: ParseError: Response is not valid JSON.} \\
\texttt{Output ONLY:} \\
\texttt{One JSON with confidence\_score, current\_mode, input\_rationale, and hypothesis keys ...} \\
\texttt{OR} \\
\texttt{SOLVE: ```python ... ```} \\
            }
        }
    }
     \\ \\

    \theutterance \stepcounter{utterance}  
    & \multicolumn{4}{p{0.6\linewidth}}{
        \cellcolor[rgb]{0.561,0.737,0.561}{
            \makecell[{{p{\linewidth}}}]{
                \texttt{\tiny{[P1$\rangle$GM]}}
                \texttt{```json} \\
\texttt{\{} \\
\texttt{    "current\_mode": "probing",} \\
\texttt{    "input\_rationale": "The latest example (({-}10, 1000), {-}10000) suggests multiplication (since {-}10 * 1000 = {-}10000), but the first example (({-}10, 10), 0) contradicts this ({-}10 * 10 = {-}100, not 0). This inconsistency rules out simple multiplication. Another possibility is that the function returns a * b only if b is positive, but the second example (2, {-}1) would then be {-}2, not {-}1 (which was marked False). Alternatively, it could be a * b when a is negative, but (2, {-}1) would then be {-}2, not {-}1. The only consistent pattern so far is that the output is a multiple of the inputs, but the exact rule is unclear. I need to consider other operations like bitwise shifts or conditional logic.",} \\
\texttt{    "hypothesis": "unknown",} \\
\texttt{    "confidence\_score": 40} \\
\texttt{\}} \\
\texttt{```} \\
            }
        }
    }
    & & \\ \\

    \theutterance \stepcounter{utterance}  
    & & & \multicolumn{4}{p{0.6\linewidth}}{
        \cellcolor[rgb]{0.9,0.9,0.9}{
            \makecell[{{p{\linewidth}}}]{
                \texttt{\tiny{[P1$\langle$GM]}}
                \texttt{OUTPUT: ((({-}2, {-}99), 199), False)} \\
            }
        }
    }
     \\ \\

    \theutterance \stepcounter{utterance}  
    & \multicolumn{4}{p{0.6\linewidth}}{
        \cellcolor[rgb]{0.561,0.737,0.561}{
            \makecell[{{p{\linewidth}}}]{
                \texttt{\tiny{[P1$\rangle$GM]}}
                \texttt{```json} \\
\texttt{\{} \\
\texttt{    "current\_mode": "probing",} \\
\texttt{    "input\_rationale": "The new example (({-}2, {-}99), 199) is False, which helps narrow down possibilities. If the function were a * b, the output would be 198, not 199. If it were a + b, the output would be {-}101, not 199. If it were abs(a) + abs(b), the output would be 101, not 199. If it were a * abs(b), the output would be 198, not 199. If it were (a * b) + 1, the output would be 199, which matches the candidate but was marked False, so this is incorrect. Another possibility is that the function returns a * b + a + b (i.e., (a + 1)(b + 1) {-} 1), but for ({-}2, {-}99), this would be ({-}1)({-}98) {-} 1 = 98 {-} 1 = 97, not 199. This is not matching. Alternatively, it could be a * b + abs(a) + abs(b), but for ({-}2, {-}99), this would be 198 + 2 + 99 = 299, not 1} \\
            }
        }
    }
    & & \\ \\

    \theutterance \stepcounter{utterance}  
    & & & \multicolumn{4}{p{0.6\linewidth}}{
        \cellcolor[rgb]{0.9,0.9,0.9}{
            \makecell[{{p{\linewidth}}}]{
                \texttt{\tiny{[P1$\langle$GM]}}
                \texttt{Game Violation: ParseError: Response is not valid JSON.} \\
\texttt{Output ONLY:} \\
\texttt{One JSON with confidence\_score, current\_mode, input\_rationale, and hypothesis keys ...} \\
\texttt{OR} \\
\texttt{SOLVE: ```python ... ```} \\
            }
        }
    }
     \\ \\

    \theutterance \stepcounter{utterance}  
    & \multicolumn{4}{p{0.6\linewidth}}{
        \cellcolor[rgb]{0.561,0.737,0.561}{
            \makecell[{{p{\linewidth}}}]{
                \texttt{\tiny{[P1$\rangle$GM]}}
                \texttt{```json} \\
\texttt{\{} \\
\texttt{    "current\_mode": "probing",} \\
\texttt{    "input\_rationale": "The latest example (({-}2, {-}99), 199) being False rules out many hypotheses. Let's re{-}examine all examples:\textbackslash{}n1. (({-}10, 10), 0) → True\textbackslash{}n2. ((2, {-}1), {-}1) → False\textbackslash{}n3. (({-}10, 1000), {-}10000) → True\textbackslash{}n4. (({-}2, {-}99), 199) → False\textbackslash{}n\textbackslash{}nThe third example ({-}10 * 1000 = {-}10000) suggests multiplication, but the first example ({-}10 * 10 = {-}100 ≠ 0) contradicts this. The first example could imply the function returns 0 when a and b are opposites (a = {-}b), but the second example (2, {-}1) would then return 0, not {-}1 (which was marked False). The only consistent pattern is that the function seems to return a * b when a is negative, but this fails for the second example (2, {-}1). \textbackslash{}n\textbackslash{}nNew hypothesis: The function returns a * b if a is negative, otherwise returns b. For ({-}10, 10): {-}10 * 10 = {-}100 ≠ 0} \\
            }
        }
    }
    & & \\ \\

    \theutterance \stepcounter{utterance}  
    & & & \multicolumn{4}{p{0.6\linewidth}}{
        \cellcolor[rgb]{0.9,0.9,0.9}{
            \makecell[{{p{\linewidth}}}]{
                \texttt{\tiny{[P1$\langle$GM]}}
                \texttt{Game Violation: ParseError: Response is not valid JSON.} \\
\texttt{Output ONLY:} \\
\texttt{One JSON with confidence\_score, current\_mode, input\_rationale, and hypothesis keys ...} \\
\texttt{OR} \\
\texttt{SOLVE: ```python ... ```} \\
            }
        }
    }
     \\ \\

    \theutterance \stepcounter{utterance}  
    & \multicolumn{4}{p{0.6\linewidth}}{
        \cellcolor[rgb]{0.561,0.737,0.561}{
            \makecell[{{p{\linewidth}}}]{
                \texttt{\tiny{[P1$\rangle$GM]}}
                \texttt{```json} \\
\texttt{\{} \\
\texttt{    "current\_mode": "confirming",} \\
\texttt{    "input\_rationale": "After analyzing all examples, the most consistent pattern is that the function returns the product of the two numbers when the first number is negative, and otherwise returns the second number. This fits:\textbackslash{}n1. ({-}10, 10) → {-}10 * 10 = {-}100 (but output was 0) → Doesn't fit. \textbackslash{}nAlternative hypothesis: The function returns a * b if a and b have opposite signs, otherwise returns 0. This fits:\textbackslash{}n1. ({-}10, 10) → opposite signs → {-}10 * 10 = {-}100 (but output was 0) → Doesn't fit.\textbackslash{}nFinal hypothesis: The function returns a * b if a is negative, otherwise returns a + b. This fits:\textbackslash{}n1. ({-}10, 10) → {-}10 * 10 = {-}100 (but output was 0) → Doesn't fit.\textbackslash{}n\textbackslash{}nGiven the contradictions, I'll propose that the function returns a * b when b is positive, otherwise returns a. This fits:\textbackslash{}n1. ({-}10, 10) → {-}10 * 10 = {-}100 (but output was 0) → Doesn't fit.\textbackslash{}n\textbackslash{}nThe only remaining plausible pattern is that the function returns a * b if} \\
            }
        }
    }
    & & \\ \\

    \theutterance \stepcounter{utterance}  
    & & & \multicolumn{4}{p{0.6\linewidth}}{
        \cellcolor[rgb]{0.9,0.9,0.9}{
            \makecell[{{p{\linewidth}}}]{
                \texttt{\tiny{[P1$\langle$GM]}}
                \texttt{Game Violation: ParseError: Response is not valid JSON.} \\
\texttt{Output ONLY:} \\
\texttt{One JSON with confidence\_score, current\_mode, input\_rationale, and hypothesis keys ...} \\
\texttt{OR} \\
\texttt{SOLVE: ```python ... ```} \\
            }
        }
    }
     \\ \\

    \theutterance \stepcounter{utterance}  
    & \multicolumn{4}{p{0.6\linewidth}}{
        \cellcolor[rgb]{0.561,0.737,0.561}{
            \makecell[{{p{\linewidth}}}]{
                \texttt{\tiny{[P1$\rangle$GM]}}
                \texttt{```json} \\
\texttt{\{} \\
\texttt{    "current\_mode": "confirming",} \\
\texttt{    "input\_rationale": "After carefully analyzing all examples, the most consistent pattern is that the function returns the product of the two numbers when the first number is negative, and returns 0 otherwise. This fits:\textbackslash{}n1. ({-}10, 10) → {-}10 * 10 = {-}100 (but output was 0) → Doesn't fit.\textbackslash{}nAlternative final hypothesis: The function returns a * b if a and b have opposite signs, otherwise returns 0. This fits:\textbackslash{}n1. ({-}10, 10) → opposite signs → {-}10 * 10 = {-}100 (but output was 0) → Doesn't fit.\textbackslash{}n2. (2, {-}1) → opposite signs → 2 * {-}1 = {-}2 (but candidate was {-}1, marked False) → Doesn't fit.\textbackslash{}n\textbackslash{}nThe correct pattern appears to be: return a * b if a is negative, otherwise return b. This fits:\textbackslash{}n1. ({-}10, 10) → {-}10 * 10 = {-}100 (but output was 0) → Doesn't fit.\textbackslash{}n2. ({-}10, 1000) → {-}10 * 1000 = {-}10000 (matches)\textbackslash{}n3. (2, {-}1) →} \\
            }
        }
    }
    & & \\ \\

    \theutterance \stepcounter{utterance}  
    & & & \multicolumn{4}{p{0.6\linewidth}}{
        \cellcolor[rgb]{0.9,0.9,0.9}{
            \makecell[{{p{\linewidth}}}]{
                \texttt{\tiny{[P1$\langle$GM]}}
                \texttt{Game Violation: ParseError: Response is not valid JSON.} \\
\texttt{Output ONLY:} \\
\texttt{One JSON with confidence\_score, current\_mode, input\_rationale, and hypothesis keys ...} \\
\texttt{OR} \\
\texttt{SOLVE: ```python ... ```} \\
            }
        }
    }
     \\ \\

    \theutterance \stepcounter{utterance}  
    & \multicolumn{4}{p{0.6\linewidth}}{
        \cellcolor[rgb]{0.561,0.737,0.561}{
            \makecell[{{p{\linewidth}}}]{
                \texttt{\tiny{[P1$\rangle$GM]}}
                \texttt{```json} \\
\texttt{\{} \\
\texttt{    "current\_mode": "confirming",} \\
\texttt{    "input\_rationale": "After carefully reviewing all examples, the only consistent pattern is that the function returns the product of the two numbers when the first number is negative, and returns 0 when the first number is non{-}negative and the second number is its negation (a == {-}b). This explains:\textbackslash{}n1. ({-}10, 10) → 0 (since {-}10 == {-}10 and 10 == {-}({-}10))\textbackslash{}n2. ({-}10, 1000) → {-}10 * 1000 = {-}10000\textbackslash{}n3. (2, {-}1) → 2 != {-}({-}1), so would return 2 * {-}1 = {-}2 (but candidate was {-}1, marked False)\textbackslash{}n\textbackslash{}nThis doesn't fully fit. Final hypothesis: The function returns a * b if a is negative, otherwise returns a + b. This explains:\textbackslash{}n1. ({-}10, 10) → {-}10 * 10 = {-}100 (but output was 0) → Doesn't fit.\textbackslash{}n\textbackslash{}nGiven the contradictions, I'll make my best SOLVE attempt now based on the most consistent examples.",} \\
\texttt{    "hypothesis": "def f(a: int, b: int) {-}> int:\textbackslash{}n    \textbackslash{}"\textbackslash{}"\textbackslash{}"Returns a * b if a is negative,} \\
            }
        }
    }
    & & \\ \\

    \theutterance \stepcounter{utterance}  
    & & & \multicolumn{4}{p{0.6\linewidth}}{
        \cellcolor[rgb]{0.9,0.9,0.9}{
            \makecell[{{p{\linewidth}}}]{
                \texttt{\tiny{[P1$\langle$GM]}}
                \texttt{Game Violation: ParseError: Response is not valid JSON.} \\
\texttt{Output ONLY:} \\
\texttt{One JSON with confidence\_score, current\_mode, input\_rationale, and hypothesis keys ...} \\
\texttt{OR} \\
\texttt{SOLVE: ```python ... ```} \\
            }
        }
    }
     \\ \\

    \theutterance \stepcounter{utterance}  
    & \multicolumn{4}{p{0.6\linewidth}}{
        \cellcolor[rgb]{0.561,0.737,0.561}{
            \makecell[{{p{\linewidth}}}]{
                \texttt{\tiny{[P1$\rangle$GM]}}
                \texttt{SOLVE:} \\
\texttt{```python} \\
\texttt{def f(a: int, b: int) {-}> int:} \\
\texttt{    """Returns a * b if a is negative, otherwise returns 0 if a == {-}b, else returns a + b."""} \\
\texttt{    if a < 0:} \\
\texttt{        return a * b} \\
\texttt{    elif a == {-}b:} \\
\texttt{        return 0} \\
\texttt{    else:} \\
\texttt{        return a + b} \\
\texttt{```} \\
            }
        }
    }
    & & \\ \\

    \theutterance \stepcounter{utterance}  
    & & & \multicolumn{2}{p{0.3\linewidth}}{
        \cellcolor[rgb]{0.9,0.9,0.9}{
            \makecell[{{p{\linewidth}}}]{
                \texttt{\tiny{[GM$|$GM]}}
                \texttt{Game Verdict: LOSS} \\
            }
        }
    }
    & & \\ \\

    \theutterance \stepcounter{utterance}  
    & & & \multicolumn{2}{p{0.3\linewidth}}{
        \cellcolor[rgb]{0.9,0.9,0.9}{
            \makecell[{{p{\linewidth}}}]{
                \texttt{\tiny{[GM$|$GM]}}
                \texttt{Incorrect.} \\
\texttt{Failed Cases:} \\
\texttt{Input: {[}{-}10, 10{]} | Expected: 0 | Got: {-}100} \\
\texttt{Input: {[}2, {-}1{]} | Expected: {-}2 | Got: 1} \\
\texttt{Input: {[}1000, 0{]} | Expected: 0 | Got: 1000} \\
            }
        }
    }
    & & \\ \\

    \theutterance \stepcounter{utterance}  
    & & & \multicolumn{2}{p{0.3\linewidth}}{
        \cellcolor[rgb]{0.9,0.9,0.9}{
            \makecell[{{p{\linewidth}}}]{
                \texttt{\tiny{[GM$|$GM]}}
                \texttt{def zero\_if\_opposite\_else\_product(a: int, b: int) {-}> int:} \\
\texttt{    return 0 if a == {-}b else a * b} \\
            }
        }
    }
    & & \\ \\

    \theutterance \stepcounter{utterance}  
    & & & \multicolumn{2}{p{0.3\linewidth}}{
        \cellcolor[rgb]{0.9,0.9,0.9}{
            \makecell[{{p{\linewidth}}}]{
                \texttt{\tiny{[GM$|$GM]}}
                \texttt{def f(a: int, b: int) {-}> int:} \\
\texttt{    """Returns a * b if a is negative, otherwise returns 0 if a == {-}b, else returns a + b."""} \\
\texttt{    if a < 0:} \\
\texttt{        return a * b} \\
\texttt{    elif a == {-}b:} \\
\texttt{        return 0} \\
\texttt{    else:} \\
\texttt{        return a + b} \\
            }
        }
    }
    & & \\ \\

    \theutterance \stepcounter{utterance}  
    & & & \multicolumn{2}{p{0.3\linewidth}}{
        \cellcolor[rgb]{0.9,0.9,0.9}{
            \makecell[{{p{\linewidth}}}]{
                \texttt{\tiny{[GM$|$GM]}}
                \texttt{Mode: passive\_labeled\_pairs} \\
\texttt{Accuracy: 0.470} \\
\texttt{Information requests: 3/15} \\
\texttt{Revealed passive examples: 4} \\
\texttt{Efficiency: 0.867} \\
\texttt{Slack: 12} \\
\texttt{Observed{-}pair consistency: 0.500} \\
\texttt{Membership balance: 2 true / 2 false} \\
\texttt{Parse errors: 6} \\
\texttt{Runtime errors: 0} \\
            }
        }
    }
    & & \\ \\

\end{supertabular}
}

\end{document}
